\section{Equip de desenvolupament}

%TODO: Definir contribucions periòdiques de cada entrega? O com o feim?

Aquest projecte és desenvolupat per un equip de quatre integrants. Com que l'objectiu principal de tots és l'aprenentatge d'una tecnologia nova (\textit{blockchain} i contractes intel·ligents), hem optat per una estructura de treball matricial. 

Tothom assumeix un rol base de \textbf{Desenvolupador Web3 amb una dedicació del 60\%}, garantint que tot l'equip participa en la creació del codi. El \textbf{40\% restant} s'assigna a un rol d'especialització per coordinar una àrea concreta, decisió basada en les fortaleses tècniques i l'experiència prèvia de cadascú.

Aquesta primera entrega ha estat elaborada de manera col·laborativa per tot l’equip, amb sessions de treball conjuntes orientades a consensuar l’abast, l’enfocament tècnic i l’organització del projecte. 
Tot i que cada membre assumeix responsabilitats principals segons el seu rol, el document s’ha revisat i validat de forma conjunta per garantir coherència i una visió compartida del sistema.

\subsection{Dylan: Responsable d'Arquitectura i \textit{Smart Contracts}}

Aquest perfil s'encarrega de dissenyar la lògica central de la xarxa, el sistema de consens i la programació dels contractes intel·ligents.

\begin{quote}
    Aporta un nivell de programació avançat i una gran capacitat de recerca autònoma. Té una visió tècnica molt àmplia i la seva principal motivació és arribar a entendre i implementar el funcionament intern de les cadenes de blocs des de zero.
\end{quote}

% \paragraph{Contribució en aquesta entrega:}
% Definició dels lliurables i recerca de tecnologies \textit{blockchain}

\subsection{Jordi: Responsable \textit{Front-end} i Experiència d'Usuari (UX/UI)}

Responsable de desenvolupar la interfície web i garantir que la interacció de l'usuari amb el sistema sigui accessible, clara i generi confiança.

\begin{quote}
    Compta amb una experiència sòlida en desenvolupament web, concretament en la creació d'interfícies d'usuari. A més, aporta una forta convicció en la necessitat de crear models descentralitzats que atorguin independència tecnològica a la societat.
\end{quote}

% \paragraph{Contribució en aquesta entrega:}
% Definició de l’abast tècnic del sistema (\textit{smart contracts}) i identificació de les dependències tècniques.

\subsection{Marc: Responsable de Qualitat (QA) i Seguretat}

Aquest perfil assegura que el codi compleixi els estàndards exigits, dissenyant les proves de \textit{software} i cercant possibles vulnerabilitats.

\begin{quote}
    Aporta experiència prèvia en la realització de proves i testatge de programari. És una persona molt curiosa que necessita reptes analítics constants per investigar on i com pot fallar un sistema.
\end{quote}

% \paragraph{Contribució en aquesta entrega:}
% Revisió de coherència funcional, definició d’evidències i traçabilitat entre requisits, proves i fites.

\subsection{Toni: Cap de Projecte}

Encarregat de supervisar la metodologia de l'equip, l'organització temporal i el compliment dels estàndards de gestió acadèmica.

\begin{quote}
    Aporta anys de bagatge professional com a desenvolupador en entorns reals i crítics (com el sector sanitari). La seva motivació és adaptar-se a tecnologies emergents i sumar la seva experiència organitzativa al desenvolupament \textit{blockchain}.
\end{quote}

% \paragraph{Contribució en aquesta entrega:} 
% Coordinació del cronograma inicial, estructuració de fites i revisió global del document.